\section{\textcolor{Salmon}{Beitrag moderner Dashboards zur Effizienzsteigerung und verbesserten Übersichtlichkeit von Legacy-Testsystemen}}\\

%Fragestellung von eigenem Problem ableiten, ohne eine Frage zu formulieren\\
%-> Subfragen bezüglich Ergebnissen, die eine Aussage zur Usability treffen\\

%kurz, nur Auflistung der Fragen\\
\noindent
\textbf{Titel}: Analyse der Steigerung der Usability durch Einführung eines Dashboards in den Legacy-Testpoolmonitoring-Prozess\\

\noindent
- aufschreiben, wie die Antworten die Thesen supporten oder nicht\\
-> die vorhergestellten Fragen\\
-> „was ist das erwartete Ergebnis“ (Regelwerk – es wird Ergebnis „turning blue“ erwartet wenn wir es richtig machen)\\
-> wie welche Antworten meine These unterstreichen\\

\noindent
- Wie passen die Hypothesen zu den Fragestellungen??\\

\noindent
- Wie ist der Stand?\\
- Worum geht es?\\
- Was ist das Ziel?\\
-> Forschungsfragen formulieren\\

\noindent
in Problemstellung:\\
- Übergang von SdF auf: was hat das mit meinem Thema zu tun\\
- wie lassen sich die Ansätze neu interpretieren auf das gegebene Problem

\noindent
Dashboard als nicht tiefe Lösung - ohne Eingriff in  das bestehende System\\
-> als erster Schritt zur Verbesserung der Usability\\
(Fazit: die die sich nicht mehr umgewöhnen wollen / können, müssen es mit diesem Ansatz nicht)

\noindent
\underline{aus Exposé}\\
Die Hauptfragestellung lautet: Welche Auswirkungen hat die Einführung eines Dashboards auf die Effizienz eines Legacy-Prozesses, untersucht am Beispiel des Testpoolmonitorings? \\
Dabei sollen folgende Teilfragen beantwortet werden:
\begin{itemize}
    \item Welche Herausforderungen bestehen, bei der Modernisierung von Legacy-Systemen?
    \item Welche Vorteile bringt ein Dashboard einem Legacy-System?
    \item Welche Rolle spielt die Akzeptanz im Nutzerumfeld, damit Modernisierungen sinnvoll sind?
\end{itemize}